\section{Discretización del problema tomográfico}

En esta sección se describe la discretización del modelo continuo presentado en el capítulo anterior, con el fin de obtener un esquema computacional para la simulación de datos tomográficos y la implementación de algoritmos de reconstrucción.

\subsection*{Discretización del dominio espacial}

Sea $f : \mathbb{R}^{2} \to \mathbb{R}$ una función con soporte compacto contenida en el cuadrado $[-1,1]^{2}$. Se considera una discretización uniforme de dicho dominio mediante una malla cartesiana de tamaño $N \times N$, dada por
\[
x_i = -1 + \frac{2(i-1)}{N-1}, \quad y_j = -1 + \frac{2(j-1)}{N-1}, \quad 1 \leq i,j \leq N.
\]

La función $f$ se aproxima entonces por una matriz $f_{i,j} \approx f(x_i,y_j)$, lo que permite interpretar la imagen como un arreglo bidimensional de valores discretos. En los experimentos numéricos se utiliza típicamente $N=256$, lo que proporciona una resolución suficiente para observar efectos de reconstrucción.

El objeto de estudio se modela mediante un \emph{phantom}, es decir, una función sintética que contiene estructuras de distintas intensidades y formas geométricas, lo que permite evaluar la capacidad de los algoritmos de reconstrucción para resolver detalles y contrastes.

\subsection*{Discretización de la transformada de Radon}

En el modelo continuo, la transformada de Radon está definida como una integral de línea sobre rectas del plano. Para su aproximación numérica, se considera un conjunto finito de ángulos
\[
\theta_k = k\,\Delta \theta, \quad k = 0,1,\dots,M-1,
\]
donde $\Delta \theta = \pi/M$. En la práctica, se trabaja con ángulos equiespaciados en el intervalo $[0,\pi)$.

Para cada ángulo $\theta_k$, la proyección de la imagen se aproxima mediante una rotación de la imagen seguida de una suma discreta en una dirección fija. Más precisamente, si $f_{\theta_k}$ denota la imagen rotada un ángulo $\theta_k$, la proyección discreta se define como
\[
g(r_i,\theta_k) \approx \sum_{j=1}^{N} f_{\theta_k}(x_i,y_j)\,\Delta r,
\]
donde $\Delta r$ es el paso de discretización en la variable espacial.

Este procedimiento proporciona una aproximación numérica de la transformada de Radon
\[
\mathcal{R}f(t,\theta),
\]
y da lugar a una matriz $g \in \mathbb{R}^{N \times M}$ que representa el \emph{sinograma} de la imagen.

\subsection*{El sinograma}

El sinograma $g(r,\theta)$ es la representación discreta del conjunto de proyecciones de la imagen. Cada columna de la matriz corresponde a una proyección en un ángulo fijo, mientras que cada fila representa la variación de la proyección respecto a la variable $r$.

Desde el punto de vista geométrico, el sinograma codifica la información de la imagen en el espacio de parámetros $(r,\theta)$, y constituye el dato de entrada del problema inverso de reconstrucción.

En los experimentos numéricos, el sinograma se calcula a partir de la discretización descrita anteriormente, y puede ser perturbado mediante la adición de ruido para simular condiciones realistas de adquisición. En particular, se consideran modelos de ruido gaussiano y de tipo Poisson, este último motivado por la naturaleza del conteo de fotones en tomografía computarizada.

\subsection*{Observaciones}

La discretización descrita anteriormente permite aproximar el operador continuo de Radon mediante un operador discreto que actúa sobre matrices. Este operador constituye la base de los algoritmos de reconstrucción considerados en esta tesis.

Cabe señalar que la aproximación mediante rotaciones y sumas discretas introduce errores numéricos asociados a la interpolación y a la discretización del dominio. Sin embargo, este enfoque es estándar en la literatura y resulta adecuado para el estudio cualitativo de los métodos de reconstrucción y del comportamiento de los filtros en el dominio de Fourier \cite{kak_slaney1988,natterer2001}.